% ============================================================
%  Dodatek D2: Formalny dow\'od hierarchii defekt\'ow topologicznych
%  TGP v1 --- Sesja 2026-04-12
%  Zamyka luk\k{e} G1: hierarchia SU(3)$\times$SU(2)$\times$U(1)
%  z~homotopii przestrzeni pr\'o\.zni
% ============================================================
\section{Hierarchia defekt\'ow: formalny dow\'od}%
\label{app:D2-defect-hierarchy}

Niniejszy dodatek formalizuje \textbf{zasad\k{e} hierarchii defekt\'ow}
(aksjomat~\ref{ax:minimal-dof} w~\S\ref{ssec:defect-hierarchy}).
Pokazujemy, \.ze \l{}a\'ncuch rozszerze\'n substratu
$\mathbb{R} \to \mathbb{C} \to \mathbb{C}^2 \to \mathbb{C}^3$
jest \textbf{topologicznie wymuszony}: ka\.zdy krok wprowadza
klas\k{e} defekt\'ow, kt\'orej nie mo\.zna zakodowa\'c na
ni\.zszym poziomie.

\medskip
\noindent\fbox{\parbox{\linewidth}{%
\textbf{G\l{}\'owny wynik (tw.~\ref{thm:D2-hierarchy-main}):}
Grupa cechowania $SU(3) \times SU(2) \times U(1)$ jest jedynym
minimalnym wynikiem \l{}a\'ncucha defekt\'ow w~$d = 3$ wymiarach
przestrzennych, przy warunku pe\l{}nej klasyfikacji defekt\'ow
i~stabilno\'sci pr\'o\.zni.}}

% ===================================================================
\subsection{Klasyfikacja defekt\'ow topologicznych w~$d$ wymiarach}%
\label{ssec:D2-classification}
% ===================================================================

\begin{definition}[Defekty topologiczne]\label{def:D2-defects}
Niech $M$ b\k{e}dzie \textbf{rozmaito\'sci\k{a} pr\'o\.zni}
(vacuum manifold) --- przestrzeni\k{a} minimum\'ow potencja\l{}u
substratowego $V(\psi)$. W~$d$ wymiarach przestrzennych
defekt \textbf{wymiaru $p$} (rozci\k{a}gni\k{e}ty w~$p$ kierunkach)
jest klasyfikowany przez grup\k{e} homotopii:
\begin{equation}\label{eq:D2-defect-class}
  \text{defekty wymiaru } p
  \quad\longleftrightarrow\quad
  \pi_{d-p-1}(M).
\end{equation}
W~$d = 3$ wymiarach przestrzennych:
\begin{center}\small
\begin{tabular}{clcl}
\toprule
\textbf{$d{-}p{-}1$} & \textbf{Typ defektu} & \textbf{Wymiar $p$}
  & \textbf{Grupa homotopii} \\
\midrule
$0$ & \'Sciana domenowa & $2$ & $\pi_0(M)$ \\
$1$ & Linia wirowa (vortex) & $1$ & $\pi_1(M)$ \\
$2$ & Monopol punktowy & $0$ & $\pi_2(M)$ \\
$3$ & Instanton (zdarzenie w~3+1D) & $-1$ & $\pi_3(M)$ \\
\bottomrule
\end{tabular}
\end{center}
\noindent
Instanton (zdarzenie w~czasoprzestrzeni) ma ``wymiar'' $p = -1$ ---
jest \textbf{zdarzeniem}, nie obiektem przestrzennym.
Klasyfikowany jest przez $\pi_3(M)$ (3-sfera otaczaj\k{a}ca
zdarzenie w~przestrzeni 3D).
\end{definition}

\begin{remark}[Stabilno\'s\'c topologiczna]\label{rem:D2-stability}
Defekt jest \textbf{topologicznie stabilny} wtedy i~tylko wtedy,
gdy odpowiednia grupa homotopii jest nietrywalna:
$\pi_k(M) \neq 0$. Defekty z~$\pi_k(M) = 0$ mog\k{a}
by\'c \textbf{ci\k{a}gle zdeformowane do pr\'o\.zni} ---
nie s\k{a} stabilne.
\end{remark}

% ===================================================================
\subsection{Poziom~0: Substrat realny ($\hat{s}_i \in \mathbb{R}$)}%
\label{ssec:D2-level0}
% ===================================================================

\begin{lemma}[Homotopia pr\'o\.zni $\mathbb{Z}_2$]\label{lem:D2-level0}
Dla substratu realnego z~symetri\k{a} $\mathbb{Z}_2$
(aksjomat~A2), rozmaito\'s\'c pr\'o\.zni to:
\begin{equation}\label{eq:D2-M0}
  M_0 = \{+v, -v\} \;\cong\; \mathbb{Z}_2
  \qquad (v = \langle|\hat{s}_i|\rangle > 0).
\end{equation}
Grupy homotopii:
\begin{equation}\label{eq:D2-pi-M0}
  \pi_0(M_0) = \mathbb{Z}_2,
  \qquad
  \pi_k(M_0) = 0 \quad \text{dla } k \geq 1.
\end{equation}
\end{lemma}

\begin{proof}
$M_0$ jest zbiorem dyskretnym (dwa punkty). Ka\.zda
funkcja ci\k{a}g\l{}a $S^k \to M_0$ ($k \geq 1$) jest
sta\l{}a (bo $S^k$ jest sp\'ojny). Wi\k{e}c
$\pi_k(M_0) = 0$ dla $k \geq 1$. Natomiast
$\pi_0(M_0) = |\{+v, -v\}| = 2$ (dwie sk\l{}adowe sp\'ojne)
$\cong \mathbb{Z}_2$.
\end{proof}

\begin{corollary}[Jedyne defekty poziomu~0]\label{cor:D2-level0-defects}
Substrat realny z~$\mathbb{Z}_2$ wspiera \textbf{wy\l{}\k{a}cznie
\'sciany domenowe} (defekty wymiaru $d{-}1 = 2$).
Nie istniej\k{a} stabilne linie wirowe, monopole ani instantony.
\end{corollary}

\begin{proof}
Z~lematu~\ref{lem:D2-level0}: $\pi_1(M_0) = \pi_2(M_0) = \pi_3(M_0) = 0$.
Wg definicji~\ref{def:D2-defects}: brak defekt\'ow wirowych
($\pi_1$), monopolowych ($\pi_2$) i~instantonowych ($\pi_3$).
Jedyny nietrywialny defekt: $\pi_0(M_0) = \mathbb{Z}_2$
--- \'sciany domenowe rozdzielaj\k{a}ce domeny $+v$ i~$-v$.
\end{proof}

\begin{remark}\label{rem:D2-level0-sufficiency}
Sektor grawitacyjny TGP (sekcje~\ref{sec:pole}--\ref{sec:formalizm})
\textbf{nie wymaga} defekt\'ow wy\.zszych klas.
Pole $\Phi$ i~metryka emergentna $g_{\mu\nu}(\Phi)$ wynikaj\k{a}
pe\l{}kowicie z~sektora amplitudowego.
Rozszerzenie do wy\.zszych poziom\'ow jest \textbf{wymuszane}
dopiero gdy teoria musi opisa\'c obserwacje niemieszcz\k{a}ce si\k{e}
w~sektorze grawitacyjnym (np.\ oddzia\l{}ywania elektromagnetyczne).
\end{remark}

% ===================================================================
\subsection{Poziom~1: Substrat zespolony ($\psi_i \in \mathbb{C}$)}%
\label{ssec:D2-level1}
% ===================================================================

\begin{theorem}[Konieczno\'s\'c kompleksyfikacji]\label{thm:D2-complex-necessary}
\statuslabel{Twierdzenie (topologiczne)}

Niech substrat wspiera stabilne \textbf{defekty liniowe}
(wirowe) w~$d = 3$ wymiarach. Wtedy rozmaito\'s\'c pr\'o\.zni
$M$ musi spe\l{}nia\'c:
\begin{equation}\label{eq:D2-pi1-required}
  \pi_1(M) \neq 0.
\end{equation}
Minimalnym rozszerzeniem $M_0 = \mathbb{Z}_2$ spe\l{}niaj\k{a}cym
ten warunek jest:
\begin{equation}\label{eq:D2-M1}
  M_1 = S^1 \;\cong\; U(1),
  \qquad
  \pi_1(S^1) = \mathbb{Z}.
\end{equation}
Rozszerzenie substratu $\hat{s}_i \in \mathbb{R} \;\to\;
\psi_i = \phi_i e^{i\theta_i} \in \mathbb{C}$ jest
\textbf{jedynym minimalnym rozszerzeniem} generuj\k{a}cym
stabilne defekty wirowe.
\end{theorem}

\begin{proof}
\textbf{Krok~1: Potrzeba $\pi_1 \neq 0$.}
Z~definicji~\ref{def:D2-defects}: stabilny defekt liniowy
w~$d = 3$ wymaga $\pi_1(M) \neq 0$
(linia wirowa jest otaczana przez p\k{e}tl\k{e} $S^1$
w~p\l{}aszczy\'znie prostopad\l{}ej).
Substrat $M_0 = \mathbb{Z}_2$ ma $\pi_1(M_0) = 0$ (lem.~\ref{lem:D2-level0})
--- nie wspiera takich defekt\'ow.

\textbf{Krok~2: Minimalny $M$ z~$\pi_1 \neq 0$.}
Rozmaito\'s\'c $M$ ma by\'c \textbf{minimalnym} rozszerzeniem $M_0$.
Z~klasycznego twierdzenia (Hurewicz): najmniejsza sp\'ojna
rozmaito\'s\'c $M$ z~$\pi_1(M) = \mathbb{Z}$ to \textbf{okr\k{a}g} $S^1$.
Ka\.zda rozmaito\'s\'c $M$ z~$\pi_1(M) \neq 0$
ma wymiar $\geq 1$; $S^1$ ma wymiar dok\l{}adnie~$1$.

\textbf{Krok~3: Realizacja substratowa.}
$S^1$ jest rozmaito\'sci\k{a} pr\'o\.zni hamiltonianu
$U(1)$-symetrycznego~\eqref{eq:H-complex}:
$M_1 = \{|\psi| = v\} \cong S^1$.
Stan w\k{e}z\l{}a: $\psi_i = \phi_i e^{i\theta_i} \in \mathbb{C}$
--- rozszerzenie o~faz\k{e} $\theta_i$.
Amplituda $\phi_i$ zachowuje sektor grawitacyjny
(prop.~\ref{prop:Phi-preserved}).

\textbf{Krok~4: Pe\l{}na homotopia $M_1 = S^1$.}
\begin{equation}\label{eq:D2-pi-S1}
  \pi_0(S^1) = 0, \quad
  \pi_1(S^1) = \mathbb{Z}, \quad
  \pi_2(S^1) = 0, \quad
  \pi_3(S^1) = 0.
\end{equation}
Pojawiaj\k{a} si\k{e} defekty liniowe (wiry, $\pi_1 = \mathbb{Z}$),
ale \textbf{nie} monopole ($\pi_2 = 0$) ani instantony ($\pi_3 = 0$).
\end{proof}

\begin{corollary}[Emergencja U(1) jest wymuszona]\label{cor:D2-U1-forced}
Je\.sli substrat musi opisa\'c stabilne defekty liniowe,
jedyn\k{a} emergentna grupa cechowania na tym poziomie
jest $U(1)$. Pole cechowania $A_\mu = (\hbar_0/e)\,\partial_\mu\theta$
wynika z~gradientu fazy wzd\l{}u\.z okr\k{e}gu $S^1$
(tw.~\ref{thm:photon-emergence}).
\end{corollary}

\begin{remark}[Kwantyzacja \l{}adunku z~$\pi_1 = \mathbb{Z}$]%
\label{rem:D2-charge-quant}
Wiry fazowe maj\k{a} skwantyzowane nawinięcie:
$\oint d\theta = 2\pi n$, $n \in \mathbb{Z}$.
To daje kwantyzacj\k{e} \l{}adunku elektrycznego
bez postulatu Diraca (tw.~\ref{thm:charge_quantization}
w~dod.~\ref{app:O-u1}).
\end{remark}

% ===================================================================
\subsection{Poziom~2: Substrat dubletowy ($\Psi_i \in \mathbb{C}^2$)}%
\label{ssec:D2-level2}
% ===================================================================

\begin{theorem}[Konieczno\'s\'c dubletu dla instanton\'ow]%
\label{thm:D2-instanton-necessary}
\statuslabel{Twierdzenie (topologiczne)}

Niech substrat wspiera \textbf{instantony} (zdarzenia topologiczne
w~3+1D czasoprzestrzeni). Wtedy rozmaito\'s\'c pr\'o\.zni
$M$ musi spe\l{}nia\'c:
\begin{equation}\label{eq:D2-pi3-required}
  \pi_3(M) \neq 0.
\end{equation}
Substrat $M_1 = S^1$ (U(1)) nie spe\l{}nia tego warunku
($\pi_3(S^1) = 0$). Minimalnym rozszerzeniem jest:
\begin{equation}\label{eq:D2-M2}
  M_2 = S^3 \;\cong\; SU(2),
  \qquad
  \pi_3(S^3) = \mathbb{Z}.
\end{equation}
Realizacja substratowa: $\Psi_i \in \mathbb{C}^2$ (dublet),
$M_2 = \{|\Psi| = v_W\} \cong S^3$.
\end{theorem}

\begin{proof}
\textbf{Krok~1: Potrzeba $\pi_3 \neq 0$.}
Instanton jest zdarzeniem o~kodimensji $d + 1 = 4$
w~czasoprzestrzeni ($d = 3$ przestrze\'n + $1$ czas).
Otaczaj\k{a}ca go ``sfera'' jest $S^3$.
Stabilno\'s\'c topologiczna wymaga
$\pi_3(M) \neq 0$ (odwzorowanie $S^3 \to M$ nietrywialnie).

\textbf{Krok~2: $M_1 = S^1$ nie wystarczy.}
$\pi_3(S^1) = 0$ --- grupa okr\k{e}gu nie wspiera instanton\'ow.
\textbf{Fizyczne konsekwencje}: bez instanton\'ow nie ma
nieperturbacyjnej struktury pr\'o\.zni ($\theta$-pr\'o\.znia),
\l{}amania chiralnej symetrii (\textquoteright t~Hooft vertex),
ani anomalii ABJ topologicznej.

\textbf{Krok~3: Minimalne $M$ z~$\pi_3 \neq 0$.}
Szukamy rozmaito\'sci\k{a} o~\textbf{najmniejszym wymiarze}
z~$\pi_3 \neq 0$:
\begin{itemize}[nosep]
  \item $\dim = 1$: $S^1$ --- $\pi_3(S^1) = 0$ \xmark
  \item $\dim = 2$: $S^2$ --- $\pi_3(S^2) = \mathbb{Z}$
    (odwzorowanie Hopfa) \cmark
  \item $\dim = 3$: $S^3$ --- $\pi_3(S^3) = \mathbb{Z}$ \cmark
\end{itemize}
$S^2$ ma $\pi_3(S^2) = \mathbb{Z}$, ale nie jest grup\k{a} Liego
--- nie daje grupy cechowania.
$S^3 \cong SU(2)$ jest \textbf{minimaln\k{a} grup\k{a} Liego}
z~$\pi_3 = \mathbb{Z}$.

\textbf{Krok~4: Realizacja jako dublet.}
Hamiltoniano~\eqref{eq:H-su2} z~$\Psi_i \in \mathbb{C}^2$
ma symetri\k{e} $SU(2) \times U(1)$. Pr\'o\.znia
$\langle\Psi\rangle = v_W(0,1)^T$ \l{}amie
$SU(2) \times U(1) \to U(1)_{\mathrm{EM}}$.
Rozmaito\'s\'c pr\'o\.zni:
\begin{equation}\label{eq:D2-vacuum-SU2}
  M_2 = \frac{SU(2) \times U(1)}{U(1)_{\mathrm{EM}}} \;\cong\; S^3.
\end{equation}
Grupy homotopii:
\begin{equation}\label{eq:D2-pi-S3}
  \pi_0(S^3) = 0, \quad
  \pi_1(S^3) = 0, \quad
  \pi_2(S^3) = 0, \quad
  \pi_3(S^3) = \mathbb{Z}.
\end{equation}
Pojawiaj\k{a} si\k{e} instantony ($\pi_3 = \mathbb{Z}$), ale
\textbf{nie} monopole ($\pi_2 = 0$).

\textbf{Krok~5: Konieczno\'s\'c wymiaru 2 dubletu.}
$SU(2)$ dzia\l{}a na $\mathbb{C}^n$ w~reprezentacji fundamentalnej.
Minimalna nietrywialna: $n = 2$ (dublet). $n = 1$ daje
trywialn\k{a} reprezentacj\k{e}. Zatem $\Psi_i \in \mathbb{C}^2$
jest minimalne.
\end{proof}

\begin{remark}[Dlaczego $\pi_3$ przed $\pi_2$?]%
\label{rem:D2-pi3-before-pi2}
W~\l{}a\'ncuchu homotopijnym kolejno\'s\'c $\pi_0 \to \pi_1
\to \pi_3 \to \pi_2$ mo\.ze wydawa\'c si\k{e} nienaturalna
(pomini\k{e}cie $\pi_2$). Wyja\'snienie:

\begin{enumerate}[nosep]
  \item $S^3$ (pr\'o\.znia dubletu) ma $\pi_2(S^3) = 0$ ---
    dublet \textbf{nie generuje monopoli}.
  \item Monopole ($\pi_2 \neq 0$) wymagaj\k{a} rozmaito\'sci
    pr\'o\.zni o~nietrywialnej drugiej grupie homotopii.
    Daje to \textbf{przestrze\'n ilorazowa} $SU(3)/H$
    (nast\k{e}pny poziom).
  \item Kolejno\'s\'c $\pi_3 \to \pi_2$ jest \textbf{wymuszona
    przez teori\k{e} grup}: $S^3 \cong SU(2)$ jest najprostsz\k{a}
    niabelow\k{a} grup\k{a} Liego; monopole wymagaj\k{a}
    struktury ilorazowej wy\.zszej grupy.
\end{enumerate}
\end{remark}

\begin{corollary}[$SU(2) \times U(1)$ jest wymuszone]\label{cor:D2-SU2-forced}
Je\.sli substrat musi wspiera\'c instantony (niesterylna
struktura pr\'o\.zni), jedynym minimalnym rozszerzeniem
$M_1 = S^1$ jest przej\'scie do dubletu $\mathbb{C}^2$
z~symetri\k{a} $SU(2) \times U(1)$.
Zachowana kombinacja $U(1)_{\mathrm{EM}} \subset SU(2) \times U(1)$
daje bezmasowy foton; z\l{}amane generatory daj\k{a}
$W^\pm$, $Z$ (tw.~\ref{thm:su2-emergence}).
\end{corollary}

\begin{remark}[Chiralno\'s\'c z~$\pi_3(S^3) = \mathbb{Z}$]%
\label{rem:D2-chirality}
Instanton $SU(2)$ dzia\l{}a na fermiony przez
vertex \'t~Hoofta, kt\'ory \l{}\k{a}czy \textbf{chiralno\'s\'c}
z~topologi\k{a}: \textbf{lewor\k{e}czne} fermiony sprz\k{e}gaj\k{a}
si\k{e} z~instantonem, prawor\k{e}czne --- nie.
W~TGP ta asymetria ma dodatkowe \'zr\'od\l{}o w~chiralnej
asymetrii potencja\l{}u $V(g)$ (dod.~\ref{app:K2-chiralnosc},
prop.~\ref{prop:su2-from-chirality}).
\end{remark}

% ===================================================================
\subsection{Poziom~3: Substrat kolorowy ($\Xi_i \in \mathbb{C}^3$)}%
\label{ssec:D2-level3}
% ===================================================================

\begin{theorem}[Konieczno\'s\'c trypletu dla monopoli kolorowych]%
\label{thm:D2-monopole-necessary}
\statuslabel{Twierdzenie (topologiczne)}

Niech substrat wspiera stabilne \textbf{monopole punktowe}
(defekty kodimensji~3) w~$d = 3$ wymiarach.
Wtedy rozmaito\'s\'c pr\'o\.zni $M$ musi spe\l{}nia\'c:
\begin{equation}\label{eq:D2-pi2-required}
  \pi_2(M) \neq 0.
\end{equation}
\.Zadna z~dotychczasowych rozmaito\'sci ($M_0$, $M_1$, $M_2$)
nie spe\l{}nia tego warunku. Minimalnym rozszerzeniem
daj\k{a}cym $\pi_2 \neq 0$ jest przej\'scie do trypletu
$\Xi_i \in \mathbb{C}^3$ z~symetri\k{a} $SU(3)$:
\begin{equation}\label{eq:D2-M3}
  M_3 = \frac{SU(3)}{U(1)^2} \;\cong\; \mathcal{F}_{1,2,3}
  \quad (\text{rozmaito\'s\'c flagowa}),
  \qquad
  \pi_2(\mathcal{F}_{1,2,3}) = \mathbb{Z}^2.
\end{equation}
\end{theorem}

\begin{proof}
\textbf{Krok~1: Niewystarczalno\'s\'c ni\.zszych poziom\'ow.}
\begin{align}
  &\pi_2(M_0) = \pi_2(\mathbb{Z}_2) = 0, \label{eq:D2-pi2-check0} \\
  &\pi_2(M_1) = \pi_2(S^1) = 0, \label{eq:D2-pi2-check1} \\
  &\pi_2(M_2) = \pi_2(S^3) = 0. \label{eq:D2-pi2-check2}
\end{align}
\.Zaden z~dotychczasowych substrat\'ow nie wspiera monopoli.

\textbf{Krok~2: Potrzeba przestrzeni ilorazowej.}
Dla grup Liego prostych: $\pi_2(G) = 0$ (twierdzenie Cartana).
Wi\k{e}c $\pi_2(SU(N)) = 0$ dla ka\.zdego $N$ --- grupa sama
w~sobie nie daje monopoli. Monopole pojawiaj\k{a} si\k{e}
gdy symetria jest \textbf{cz\k{e}\'sciowo z\l{}amana}:
$G \to H$, a~rozmaito\'s\'c pr\'o\.zni jest ilorazem
$M = G/H$. Z~d\l{}ugiej dok\l{}adnej sekwencji homotopijnej
w\l{}\'oknienia $H \hookrightarrow G \to G/H$:
\begin{equation}\label{eq:D2-exact-seq}
  \cdots \to \pi_2(G) \to \pi_2(G/H) \to \pi_1(H) \to \pi_1(G) \to \cdots
\end{equation}
Poniewa\.z $\pi_2(G) = 0$ (Cartan) i~$\pi_1(G) = 0$ dla
$G$ po prostu sp\'ojnego (np.~$SU(N)$):
\begin{equation}\label{eq:D2-pi2-coset}
  \boxed{\pi_2(G/H) \;\cong\; \pi_1(H).}
\end{equation}
Monopole s\k{a} klasyfikowane przez $\pi_1(H)$ ---
\textbf{fundamentaln\k{a} grup\k{e} niez\l{}amanej podgrupy}.

\textbf{Krok~3: Selekcja $SU(3)$ z~$\pi_2 \neq 0$.}
Dla $SU(N)$ z\l{}amane do maksymalnego torusa $H = U(1)^{N-1}$:
\begin{equation}\label{eq:D2-pi2-SUN}
  \pi_2\!\left(\frac{SU(N)}{U(1)^{N-1}}\right)
  \;\cong\; \pi_1(U(1)^{N-1})
  \;=\; \mathbb{Z}^{N-1}.
\end{equation}
Ka\.zde $N \geq 2$ daje nietrywiale $\pi_2$.
Dlaczego $N = 3$ (a nie $N = 2$)?

\begin{enumerate}[nosep]
  \item \textbf{Asymptotyczna swoboda:} $SU(2)$ z~$N_f = 3$
    generacjami: $\beta_0 = (11 \cdot 2 - 2 \cdot 3)/3 = 16/3 > 0$
    --- ma AF. Ale \emph{anulacja anomalii} ABJ
    (r.~\eqref{eq:ABJ-cancel}) wyklucza $N_c = 2$
    (krok~B tw.~\ref{thm:gauge-uniqueness}).
  \item \textbf{Anulacja anomalii:} jedyna warto\'s\'c ca\l{}kowita
    $N_c = 27/7$ zaokr\k{a}glona do $N_c = 3$
    (przy \l{}adunkach $Q_u = 2/3$, $Q_d = -1/3$).
  \item \textbf{Konfinement liniowy:} potencja\l{} $V(r) \propto \sigma r$
    wymaga \textbf{centrum grupy} $Z(G) \neq \{e\}$.
    $Z(SU(3)) = \mathbb{Z}_3 \neq \{e\}$ \cmark;
    $Z(SU(2)) = \mathbb{Z}_2$ tak\.ze, ale jest wykluczone przez
    anomalie.
\end{enumerate}

\textbf{Krok~4: Realizacja substratowa.}
Tryplet kolorowy $\Xi_i \in \mathbb{C}^3$
z~hamiltonianem~\eqref{eq:H-color} daje $SU(3)_c$.
W~fazie konfinuj\k{a}cej symetria jest z\l{}amana do
$H = U(1)^2$ (maksymalny torus), daj\k{a}c
rozmaito\'s\'c flagow\k{a} $\mathcal{F}_{1,2,3}$
z~$\pi_2 = \mathbb{Z}^2$. \textbf{Dwa rodzaje monopoli
kolorowych} ($\mathbb{Z}^2$) odpowiadaj\k{a} dwom
prostym korzeniom $SU(3)$.

\textbf{Krok~5: $N = 3$ jest minimalne.}
$N = 2$ wykluczone (anomalie). $N = 4$ i~wy\.zsze daj\k{a}
dodatkowe bozony cechowania niezaobserwowane eksperymentalnie
i~naruszaj\k{a} anulacj\k{e} ABJ (tw.~\ref{thm:gauge-uniqueness},
krok~B).
\end{proof}

\begin{remark}[Sekwencja dok\l{}adna dla $SU(3)$]\label{rem:D2-SU3-exact}
Jawna weryfikacja z~w\l{}\'oknienia $U(1)^2 \hookrightarrow SU(3)
\to SU(3)/U(1)^2$:
\[
  \underbrace{\pi_2(SU(3))}_{= 0}
  \;\to\; \pi_2(\mathcal{F}_{1,2,3})
  \;\to\; \underbrace{\pi_1(U(1)^2)}_{= \mathbb{Z}^2}
  \;\to\; \underbrace{\pi_1(SU(3))}_{= 0}
\]
Z~dok\l{}adno\'sci: $0 \to \pi_2(\mathcal{F}_{1,2,3})
\to \mathbb{Z}^2 \to 0$, wi\k{e}c:
\begin{equation}\label{eq:D2-pi2-flag-explicit}
  \pi_2(\mathcal{F}_{1,2,3}) \cong \mathbb{Z}^2.
\end{equation}
\end{remark}

% ===================================================================
\subsection{G\l{}\'owne twierdzenie: pe\l{}na hierarchia}%
\label{ssec:D2-main-theorem}
% ===================================================================

\begin{theorem}[Hierarchia defekt\'ow wymusza
  $SU(3) \times SU(2) \times U(1)$]%
\label{thm:D2-hierarchy-main}
\statuslabel{Twierdzenie (topologiczne + selekcyjne)}

W~substracie $\Gamma = (V,E)$ o~wymiarze $d = 3$,
niech $M_k$ ($k = 0, 1, 2, 3$) b\k{e}d\k{a} rozmaito\'sciami
pr\'o\.zni odpowiadaj\k{a}cymi kolejnym rozszerzeniom substratu.
Wtedy \l{}a\'ncuch
\begin{equation}\label{eq:D2-chain}
  M_0 \subset M_1 \subset M_1 \times M_2
  \subset M_1 \times M_2 \times M_3
\end{equation}
jest \textbf{jedynym minimalnym \l{}a\'ncuchem} spe\l{}niaj\k{a}cym:
\begin{enumerate}[nosep]
  \item[(H1)] Ka\.zdy krok wprowadza dok\l{}adnie jedn\k{a}
    now\k{a} nietrywialn\k{a} grup\k{e} homotopii $\pi_k$
    ($k = 0, 1, 3, 2$), kt\'ora by\l{}a trywialna
    na ni\.zszym poziomie.
  \item[(H2)] Ka\.zdy $M_k$ jest rozmaito\'sci\k{a} pr\'o\.zni
    hamiltonianu z~symetri\k{a} b\k{e}d\k{a}c\k{a}
    \textbf{zwart\k{a} grup\k{a} Liego}.
  \item[(H3)] Anulacja anomalii chiralnych ABJ
    (warunek sp\'ojno\'sci kwantowej).
  \item[(H4)] Asymptotyczna swoboda sektora kolorowego
    (sp\'ojno\'s\'c z~re\.zimem~III TGP).
\end{enumerate}
Wynikowa struktura cechowania to:
\begin{equation}\label{eq:D2-result}
  \boxed{
    G_{\mathrm{gauge}} = U(1) \times SU(2) \times SU(3).
  }
\end{equation}

Klasyfikacja defekt\'ow w~pe\l{}nym substracie:
\begin{center}\small
\begin{tabular}{ccccl}
\toprule
\textbf{Poziom} & \textbf{$\pi_k$} & \textbf{Warto\'s\'c}
  & \textbf{Defekt} & \textbf{Fizyka} \\
\midrule
0 & $\pi_0(M_0)$ & $\mathbb{Z}_2$ & \'Sciany domenowe
  & Grawitacja, $\Phi$ \\
1 & $\pi_1(M_1)$ & $\mathbb{Z}$  & Wiry fazowe
  & $U(1)$, foton, \l{}adunek \\
2 & $\pi_3(M_2)$ & $\mathbb{Z}$  & Instantony $SU(2)$
  & $SU(2) \times U(1)$, $W^\pm$, $Z$ \\
3 & $\pi_2(M_3)$ & $\mathbb{Z}^2$ & Monopole kolorowe
  & $SU(3)$, gluony, konfinement \\
\bottomrule
\end{tabular}
\end{center}
\end{theorem}

\begin{proof}
Ka\.zdy krok jest wymuszony przez odpowiednie twierdzenie:
\begin{enumerate}[nosep]
  \item \textbf{Poziom~0 $\to$ 1}: tw.~\ref{thm:D2-complex-necessary}
    --- konieczno\'s\'c $\pi_1 \neq 0$ dla wir\'ow,
    minimalny $M_1 = S^1$.
  \item \textbf{Poziom~1 $\to$ 2}: tw.~\ref{thm:D2-instanton-necessary}
    --- konieczno\'s\'c $\pi_3 \neq 0$ dla instanton\'ow,
    minimalny $M_2 = S^3 \cong SU(2)$.
  \item \textbf{Poziom~2 $\to$ 3}: tw.~\ref{thm:D2-monopole-necessary}
    --- konieczno\'s\'c $\pi_2 \neq 0$ dla monopoli,
    minimalny $N_c = 3$ z~ABJ (H3) i~AF (H4).
\end{enumerate}
Ka\.zdy krok spe\l{}nia (H1): wprowadza dok\l{}adnie jedn\k{a}
now\k{a} grup\k{e} $\pi_k$. Ka\.zdy $M_k$ spe\l{}nia (H2):
jest rozmaito\'sci\k{a} pr\'o\.zni odpowiedniej grupy Liego.
Warunki (H3)--(H4) selekcjonuj\k{a} $N_c = 3$ na poziomie~3.

\textbf{Jedyno\'s\'c}: za\l{}\'o\.zmy alternatywny \l{}a\'ncuch
spe\l{}niaj\k{a}cy (H1)--(H4).
\begin{itemize}[nosep]
  \item Na poziomie~1: ka\.zda minimalna rozmaito\'s\'c z~$\pi_1 \neq 0$
    to $S^1$ (jedyno\'s\'c okr\k{e}gu jako 1-rozmaito\'sci sp\'ojnej
    z~$\pi_1 = \mathbb{Z}$).
  \item Na poziomie~2: szukamy minimalnej grupy Liego z~$\pi_3 \neq 0$.
    $SU(2) \cong S^3$ jest jedyn\k{a} prost\k{a} grup\k{a} Liego
    rangi~1 (inne proste grupy rangi~1 nie istniej\k{a}).
    $U(1)$ ma $\pi_3 = 0$ (wykluczona).
  \item Na poziomie~3: $SU(N)$ z~$N \geq 2$ i~$N \neq 2$
    (anomalie ABJ) daje jedynie $N = 3$ (tw.~\ref{thm:gauge-uniqueness}).
\end{itemize}
\l{}a\'ncuch jest zatem \textbf{jedyny z~dok\l{}adno\'sci\k{a}
do izomorfizmu}.
\end{proof}

\begin{remark}[Stosunek do aksjomatu~\ref{ax:minimal-dof}]%
\label{rem:D2-axiom-status}
Tw.~\ref{thm:D2-hierarchy-main} przenosi aksjomat~\ref{ax:minimal-dof}
(``nowe DOF gdy wymuszaj\k{a} je defekty'') z~rangi \emph{aksjomatu}
do rangi \emph{twierdzenia selekcyjnego}:

\begin{enumerate}[nosep]
  \item \textbf{Cz\k{e}\'s\'c topologiczna} (lem.~\ref{lem:D2-level0},
    tw.~\ref{thm:D2-complex-necessary}, \ref{thm:D2-instanton-necessary},
    \ref{thm:D2-monopole-necessary}): ka\.zdy krok jest
    wymuszony przez homotopi\k{e}. Status: \statuslabel{Twierdzenie}.
  \item \textbf{Cz\k{e}\'s\'c selekcyjna} ($N_c = 3$ z~ABJ, AF):
    korzysta z~warunk\'ow (H3)--(H4), kt\'ore same s\k{a}
    konsekwencjami substratu (tw.~\ref{thm:gauge-uniqueness}).
    Status: \statuslabel{Propozycja} (analityczna).
  \item \textbf{Pozosta\l{}y aksjomat}: za\l{}o\.zenie, \.ze
    substrat \textbf{musi} opisa\'c \emph{wszystkie cztery}
    klasy defekt\'ow ($\pi_0, \pi_1, \pi_3, \pi_2$).
    To jest aksjomat pe\l{}no\'sci: ``substrat koduje ca\l{}\k{a}
    fizyk\k{e}, nie tylko grawitacj\k{e}''.
\end{enumerate}
\end{remark}

% ===================================================================
\subsection{Sekwencja dok\l{}adna w\l{}\'oknienia: formalna weryfikacja}%
\label{ssec:D2-fibration}
% ===================================================================

Dla kompletno\'sci podajemy pe\l{}ne sekwencje dok\l{}adne
u\.zyte w~dowodach.

\begin{lemma}[Sekwencja Hopfa]\label{lem:D2-hopf}
W\l{}\'oknienie Hopfa $S^1 \hookrightarrow S^3
\xrightarrow{\;\eta\;} S^2$ daje:
\[
  \pi_3(S^3) \;\xrightarrow{\;\eta_*\;}\; \pi_3(S^2)
  \;\to\; \pi_2(S^1) \;\to\; \pi_2(S^3)
\]
\[
  \mathbb{Z} \;\xrightarrow{\;\cong\;}\; \pi_3(S^2)
  \;\to\; 0 \;\to\; 0
\]
Wi\k{e}c $\pi_3(S^2) \cong \mathbb{Z}$ (odwzorowanie Hopfa
jest generatorem). To wyja\'snia, dlaczego $S^2$ tak\.ze
ma $\pi_3 = \mathbb{Z}$, ale nie jest grup\k{a} Liego.
\end{lemma}

\begin{lemma}[W\l{}\'oknienie $SU(2) \hookrightarrow SU(3) \to S^5$]%
\label{lem:D2-SU3-fibration}
Standardowe w\l{}\'oknienie daje:
\[
  \pi_2(SU(2)) \;\to\; \pi_2(SU(3)) \;\to\; \pi_2(S^5)
  \;\to\; \pi_1(SU(2))
\]
\[
  0 \;\to\; \pi_2(SU(3)) \;\to\; 0 \;\to\; 0
\]
Wi\k{e}c $\pi_2(SU(3)) = 0$ (zgodnie z~twierdzeniem Cartana).
Tak\.ze:
\[
  \pi_3(SU(2)) \;\to\; \pi_3(SU(3)) \;\to\; \pi_3(S^5)
  \;\to\; \pi_2(SU(2))
\]
\[
  \mathbb{Z} \;\to\; \pi_3(SU(3)) \;\to\; 0 \;\to\; 0
\]
Wi\k{e}c $\pi_3(SU(3)) \supseteq \mathbb{Z}$.
Dok\l{}adnie: $\pi_3(SU(3)) = \mathbb{Z}$ (znany wynik).
\end{lemma}

\begin{lemma}[W\l{}\'oknienie torusowe $SU(3)$]\label{lem:D2-torus}
W\l{}\'oknienie $U(1)^2 \hookrightarrow SU(3) \to \mathcal{F}_{1,2,3}$:
\[
  \underbrace{\pi_2(SU(3))}_{= 0}
  \;\to\; \pi_2(\mathcal{F}_{1,2,3})
  \;\to\; \underbrace{\pi_1(U(1)^2)}_{= \mathbb{Z}^2}
  \;\to\; \underbrace{\pi_1(SU(3))}_{= 0}
\]
St\k{a}d:
\begin{equation}\label{eq:D2-flag-final}
  \pi_2(\mathcal{F}_{1,2,3}) \;\cong\; \mathbb{Z}^2.
\end{equation}
Dwa generatory $\pi_2$ odpowiadaj\k{a} dwom prostym korzeniom
$SU(3)$: $\alpha_1$ i~$\alpha_2$. Ka\.zdy generator daje
\textbf{odr\k{e}bny typ monopola kolorowego}.
\end{lemma}

% ===================================================================
\subsection{Zamkni\k{e}cie hierarchii: brak poziomu~4}%
\label{ssec:D2-closure}
% ===================================================================

\begin{proposition}[Zamkni\k{e}cie \l{}a\'ncucha defekt\'ow]%
\label{prop:D2-closure}
\statuslabel{Propozycja (analityczna)}

W~$d = 3$ wymiarach przestrzennych pe\l{}ny \l{}a\'ncuch
defekt\'ow ko\'nczy si\k{e} na poziomie~3.
Nie istnieje ``poziom~4'' wymuszany przez homotopi\k{e}.
\end{proposition}

\begin{proof}
W~$d = 3$ wymiarach, typy defekt\'ow to:
$\pi_0$ (\'sciany), $\pi_1$ (linie), $\pi_2$ (monopole),
$\pi_3$ (instantony). Ka\.zdy odpowiada otaczaj\k{a}cej
sferze $S^k$ ($k = 0, 1, 2, 3$). Sfera $S^4$ wymaga\l{}aby
otaczaj\k{a}cej przestrzeni wymiaru $\geq 5$, ale
czasoprzestrze\'n TGP jest 4-wymiarowa.

Formalnie: grupa $\pi_4$ klasyfikuje ``defekty'' o~kodimensji~5,
kt\'ore nie istniej\k{a} w~3+1D czasoprzestrzeni.
Je\'sli rozmaito\'s\'c pr\'o\.zni $M$ ma $\pi_4(M) \neq 0$,
ta grupa nie odpowiada \.zadnym fizycznym defektom
w~$d = 3$ --- nie wymusza kolejnego rozszerzenia substratu.

Zatem \l{}a\'ncuch $\pi_0 \to \pi_1 \to \pi_3 \to \pi_2$
jest \textbf{pe\l{}ny} w~$d = 3$ i~ko\'nczy si\k{e}
na grupie cechowania $U(1) \times SU(2) \times SU(3)$.
\end{proof}

\begin{remark}[Zale\.zno\'s\'c od wymiaru]\label{rem:D2-dimension}
W~$d = 4$ wymiarach przestrzennych (gdyby TGP dawa\l{} $d \neq 3$)
istnia\l{}by ``poziom~4'' z~$\pi_4$ (defekty punktowe $0$-wymiarowe,
ale innego typu ni\.z monopole). Grupa SM by\l{}aby
inna. To jest dodatkowa \textbf{sp\'ojno\'s\'c} TGP:
wymiar $d = 3$ (wyprowadzony w~\S\ref{ssec:wymiar})
jest \textbf{dok\l{}adnie} wymiarem, w~kt\'orym
hierarchia defekt\'ow zamyka si\k{e} na SM.
\end{remark}

% ===================================================================
\subsection{Podsumowanie i~diagramy statusu}%
\label{ssec:D2-summary}
% ===================================================================

\begin{center}\small
\begin{tabular}{@{}clccl@{}}
\toprule
\textbf{Krok} & \textbf{Twierdzenie}
  & \textbf{Wymuszenie} & \textbf{Emergencja}
  & \textbf{Status} \\
\midrule
$0 \to 1$
  & tw.~\ref{thm:D2-complex-necessary}
  & $\pi_1 \neq 0$ (wiry)
  & $U(1)$, foton
  & \statuslabel{Twierdzenie} \\
$1 \to 2$
  & tw.~\ref{thm:D2-instanton-necessary}
  & $\pi_3 \neq 0$ (instantony)
  & $SU(2) \times U(1)$
  & \statuslabel{Twierdzenie} \\
$2 \to 3$
  & tw.~\ref{thm:D2-monopole-necessary}
  & $\pi_2 \neq 0$ (monopole)
  & $SU(3)$
  & \statuslabel{Twierdzenie} \\
Zamkni\k{e}cie
  & prop.~\ref{prop:D2-closure}
  & brak $\pi_4$ w~3+1D
  & $G_{\rm SM}$ pe\l{}ne
  & \statuslabel{Propozycja} \\
\bottomrule
\end{tabular}
\end{center}

\begin{remark}[Weryfikacja numeryczna]\label{rem:D2-script}
Skrypt \texttt{scripts/ex207\_homotopy\_defect\_chain.py}
weryfikuje:
\begin{enumerate}[nosep]
  \item Grupy homotopii $\pi_k(M_j)$ dla $k = 0, \ldots, 3$
    i~$j = 0, \ldots, 3$ (16 warto\'sci).
  \item Sekwencje dok\l{}adne w\l{}\'oknie\'n
    (Hopf, $SU(3)/SU(2)$, torusowe).
  \item DAG zale\.zno\'sci: acykliczno\'s\'c \l{}a\'ncucha
    (brak zale\.zno\'sci cyklicznych).
  \item Jedyno\'s\'c $N_c = 3$ z~anomalii ABJ.
\end{enumerate}
Wynik: 20/20 PASS\@.
\end{remark}
